

\section{Interactivity Evaluation}
\label{sec:evaluation}
Despite progress, the prevailing evaluation benchmarks~\cite{zhang2021flow, xie2022vfhq, zhu2022celebv} for single-person talking head generation are inadequate for assessing natural interactions among characters.
Although InterActHuman~\cite{wang2025interacthuman} introduces a comparable benchmark, its test set is limited to scenarios with only one speaker, which is not conducive to evaluating interactions among multiple characters.
To fill this gap, we have curated a collection of videos featuring two distinct identities, sourced from the web for evaluation purposes.

\subsection{Dataset Construction}
We select interactive two-person videos to construct the video dataset, named \textit{InteractiveEyes}.
Two clips of these videos are illustrated in~\cref{fig:interactivity}. 
Each video is approximately 10 seconds in duration and showcases exactly two faces throughout the entire segment. 
Furthermore, through a meticulous manual process, we segment the audio of each video to ensure that the majority of the videos capture scenes of both individuals engaging in \textit{speaking} and \textit{listening}, as well as a variety of rich eye interaction scenarios, as shown in~\cref{fig:time}. 
We have also ensured that each video includes instances of mutual gaze and head movements to provide authentic references.

\begin{figure}[t]
    \centering
\includegraphics[width=0.7\linewidth]{fig/time.pdf}
    \caption{
    Listening and speaking periods of each speaker.
    }
    \label{fig:time}
\end{figure}

\subsection{Proposed Interactivity Metric}
In addition to this dataset, we introduce a novel metric, the eye-focused \textit{Interactivity}, designed to assess the natural interaction between speakers and listeners. 
Since eye interaction is a fundamental and spontaneous behavior in conversational contexts, we use it as a key indicator of interactivity.
Drawing inspiration from the Hand Keypoint Variance (HKV) metric employed in CyberHost~\cite{lin2025cyberhost}, we propose a quantitative evaluation of the interaction by tracking the motion amplitude of eye keypoints.

To achieve this, we define $Motion$ on the sequence of face-aligned eye keypoints extracted from the generated frames, where $S$ denotes the frame sequence and $E$ the eye keypoints.
The $Motion$ is calculated as follows
\begin{equation}
    Motion = \frac{1}{|S| - 1} \sum_{j=1}^{|S| - 1} \left( \frac{1}{|E|} \sum_{i=1}^{|E|} \vert E_{i,j+1} - E_{i,j}\vert \right).
\end{equation}
Here, $i$ and $j$ denote the eye keypoint index and the frame index, while \( E_{i,j} \) denotes eye keypoints present in each frame.
This formula intuitively computes the displacement and rotation of the eye region.
We then calculate the motion during listening periods. 
The reason is, most generation methods perform well when activating the speaking subject, but the listening subject often appears rigid. 
Therefore, evaluating during listening periods is more targeted and valuable.
The lengths of the listening and speaking periods of each person are described in~\cref{fig:time}, denoted as \( L_1, L_2, L_3, L_4 \), respectively.
To quantify the responsiveness of the generated avatars, we compute the average motion intensity during the listening phases $L_2$ and $L_3$:
\begin{equation}
\begin{aligned}
\text{Interactivity} &= \frac{L_2 \cdot Motion_{L2} + L_3 \cdot Motion_{L3}}{L_2 + L_3}.
\end{aligned}
\end{equation}
This metric measures the interactivity effectively in the generated multi-character videos. 
% As the experimental results suggest, the proposed metric aligns well with human perception.
As~\cref{fig:interactivity} suggests, the proposed metric aligns well with human perception: static or sluggish eye movements receive low $Motion$ scores, while head turns and eyebrow raises increase the score thus indicating higher interactivity.
Moreover, to avoid misjudging abnormal eye movements, we implemented an exclusion algorithm detailed in the supplementary materials.


\begin{table*}[t]
    \centering
    \small
    \caption{Quantitative comparison with other competing methods on HDTF~\cite{zhang2021flow} and VFHQ~\cite{xie2022vfhq} benchmark.
    Here, OmniHuman-1.5$^{*}$~\cite{jiang2025omnihuman} refers to its ``Master Mode'' version accessed via the JiMeng platform~\cite{jimeng_platform}, which currently does not support multi-person generation.
    }
    \begin{tabular}{ccccccccccc}
        \toprule
        \multirow{2}{*}{\textbf{Method}} & \multicolumn{2}{c}{\textbf{Supported Generation Scope}} & \multicolumn{4}{c}{\textbf{HDTF}} & \multicolumn{4}{c}{\textbf{VFHQ}} \\
        \cmidrule(lr){2-3} \cmidrule(lr){4-7} \cmidrule(lr){8-11}
        & \textbf{multi-person} & \textbf{body} & \textbf{Sync-C$\uparrow$} & \textbf{FID} & \textbf{FVD} & \textbf{ID$\uparrow$} & \textbf{Sync-C$\uparrow$} & \textbf{FID} & \textbf{FVD} & \textbf{ID$\uparrow$} \\
        \midrule
        AniPortrait~\cite{wei2024aniportrait} &  $\times$ & $\times$ & 3.44 & 18.74 & 241.84 & \underline{0.94} & 2.63 & 28.54 & 269.24 & \textbf{0.95} \\
        FantasyTalking~\cite{wang2025fantasytalking} &  $\times$ & $\checkmark$ & 3.97 & 14.93 & 166.79 & 0.93 & 3.57 & 24.83 & 272.13 & \underline{0.94} \\
        StableAvatar~\cite{tu2025stableavatar} &  $\times$ & $\checkmark$  & 4.11 & 14.67 & 166.44 & 0.91 & 3.53 & 22.91 & 275.73 & 0.88 \\
        EchoMimic~\cite{chen2025echomimic} &  $\times$ & $\times$ & 5.23 & 61.53 & 381.55 & \textbf{0.95} & 4.87 & 58.72 & 486.75 & 0.86 \\
        Hallo3~\cite{cui2025hallo3} & $\times$ & $\times$ & 7.53 & 17.12 & 195.61 & 0.91 & 6.32 & 41.26 & 371.24 & 0.91 \\
        Sonic~\cite{ji2025sonic} &  $\times$ & $\times$ & 7.81 & 52.96 & 286.12 & \textbf{0.95} & 7.71 & 36.68 & 385.37 & 0.89 \\
        OmniHuman-1.5$^{*}$~\cite{jiang2025omnihuman} &  $\times$ & $\checkmark$   & 7.23  & 35.26  & 173.23  & 0.90  & 7.67  & 35.36  & 283.39  & 0.90  \\
        MuitiTalk~\cite{kong2025let} &  $\checkmark$ & $\checkmark$ & \underline{8.91} & \textbf{13.54} & \underline{162.58} & 0.93 &  \underline{7.77} & 24.25 & \textbf{243.66} & \underline{0.94} \\
        \midrule
        AnyTalker-1.3B & $\checkmark$ & $\checkmark$ & 6.85 & 14.47 & 218.01 & 0.91 & 5.81 & \underline{21.88} & \underline{267.08} & 0.91 \\
        AnyTalker-14B & $\checkmark$ & $\checkmark$ & \textbf{9.05} & \underline{13.84} & \textbf{160.87} & \underline{0.94} & \textbf{7.79} & \textbf{20.99} & 290.73 & \underline{0.94} \\
        \bottomrule
    \end{tabular}
    \label{tab:single-benchmark}
\end{table*}

\begin{table}[t]
    \centering
    \small
    \caption{
        Quantitative comparison with other competing methods on the multi-person benchmark, InteractiveEyes.
    }
    \begin{tabular}{ccccc}
    \toprule
    \textbf{Method} & \textbf{Interactivity$\uparrow$} & \textbf{Sync-C$^{*}$$\uparrow$} & \textbf{FVD} \\
    \midrule
    Ground Truth & 0.77 & 6.01 & 0 \\
     \midrule
    Bind-Your-Avatar~\cite{huang2025bind} & 0.45 & 3.03 & 695.58  \\
    MuitiTalk~\cite{kong2025let} & 0.49 & \underline{6.88} & 500.03  \\
    AnyTalker-1.3B & \underline{0.97} & 4.56 & \underline{467.84}  \\
    AnyTalker-14B & \textbf{1.01} & \textbf{6.99} & \textbf{424.15}   \\
    \bottomrule
    \end{tabular}
    \label{tab:multi-benchmark}
\end{table}